%% kapa: exam=1 round=36 year=2025 period=1 form=A subject=civil_law subject_ko=민법 seq=17 number=17
%% source: Q-Net 감정평가사 기출문제 (cst00309) artlSeq=5231525 file=raw/1008/5231525/2025년 제36회 감정평가사 1차시험 1교시 문제지 원본.pdf
%% conversion: kordoc 4.24.51 -> tools/gen1.py (auto); review=proofread-sonnet; flags=none
\begin{question}{17}
조건과 기한에 관한 설명으로 옳지 않은 것은? (다툼이 있으면 판례에 따름)
\begin{choices}
  \item 조건을 붙이는 것이 허용되지 않는 법률행위에 조건을 붙인 경우 그 조건만 무효로 된다.
  \item 법률행위의 부관으로서의 조건이 되기 위해서는 법률행위에 조건을 붙이고자 하는 의사와 그 표시가 필요하다.
  \item 기한은 채무자의 이익을 위한 것으로 추정한다.
  \item 건축허가를 받지 못하면 무효로 한다는 약정 아래 이루어진 토지매매계약은 해제조건부 계약이다.
  \item 당사자가 불확정한 사실이 발생한 때를 이행기한으로 정한 경우, 그 사실의 발생이 불가능하게 된 때에는 그 이행기한은 도래한 것으로 보아야 한다.
\end{choices}
\end{question}
\answer{1}
